\section*{Chapter \ref{ch:nw:on-chain-connectivity} --- On-Chain Connectivity}

\begin{solution}{exo:nw:on-chain-connectivity:1}
$4 \times 300 = 1\,200$~ms, 1.2 seconds; but the effective slot time is changed dynamically, so read it from the chain rather than assume the default.
\end{solution}

\begin{solution}{exo:nw:on-chain-connectivity:2}
A propagation network carries the transaction faster and with less variance than public gossip, and may still deliver it to the public; a private relay
keeps it out of the public mempool until a builder includes it. The first sells latency, the second invisibility.
\end{solution}

\begin{solution}{exo:nw:on-chain-connectivity:3}
Its peers are random: a node in the same city is usually reached through other nodes, each adding its processing, and in the model the median node in Tokyo is
reached after eight hops, \qty{40}{\milli\second} against \qty{1}{\milli\second} through the Tokyo engine.
\end{solution}

\begin{solution}{exo:nw:on-chain-connectivity:4}
Square-root fan-out replaces some pushes by an announcement and a pull, which lengthens those links by a round trip; by the proposition, longer links
cannot shorten any shortest path. It is used because it sends each transaction in full to far fewer peers, saving bandwidth for a small cost in time.
\end{solution}

\begin{solution}{exo:nw:on-chain-connectivity:5}
The Singapore engine: \qty{39.9}{\milli\second} against \qty{58.9}{\milli\second} by gossip, \qty{19}{\milli\second} sooner.
\end{solution}

\begin{solution}{exo:nw:on-chain-connectivity:6}
A tenth of a second of propagation is under one percent of a 12-second slot: nearly every transaction sent early in the slot arrives in time wherever the
proposer is, and geography matters only for what is sent at the slot's end.
\end{solution}

\begin{solution}{exo:nw:on-chain-connectivity:7}
Half the nodes after \qty{80.3}{\milli\second} and nine-tenths after \qty{89.5}{\milli\second}: 4.7 and \qty{10.0}{\milli\second} faster, because more
peers mean fewer hops.
\end{solution}

\begin{solution}{exo:nw:on-chain-connectivity:8}
The fastest path is to the next leader, not to the firm: an engine next to the firm's servers helps only when the leader is also nearby. With leaders on
several continents, the endpoint must be chosen per leader from the schedule.
\end{solution}

\begin{solution}{pb:nw:on-chain-connectivity:1}
\textbf{Part I.}
\begin{enumerate}
\item \qty{85.0}{\milli\second} and \qty{99.5}{\milli\second}.
\item \qty{113.9}{\milli\second} and \qty{149.0}{\milli\second}.
\item After \qty{83.5}{\milli\second}.
\item It takes three hops, each with \qty{5}{\milli\second} of processing, where the direct path has one engine at \qty{1}{\milli\second}.
\end{enumerate}
\textbf{Part II.}
\begin{enumerate}[resume]
\item Frankfurt's; \qty{69.5}{\milli\second}.
\item \qty{45.6}{\milli\second} one way; the model multiplies it by 1.5.
\item \qty{1}{\milli\second}, an assumption; the rest is the route.
\item The transaction itself: gossip shows it to every node on the way, the engine only to itself and the leader.
\end{enumerate}
\textbf{Part III.}
\begin{enumerate}[resume]
\item 61\% with the engine, 48\% by gossip.
\item 80.5\% against 74.0\%.
\item Once the time left exceeds the slowest arrival (\qty{99.5}{\milli\second} by gossip), both methods reach every leader; the fixed difference in
  milliseconds is then a smaller part of the window.
\item The same leader holds four slots, so the endpoint changes only at the handovers; transactions sent near the end of a leader's run go to the next
  leader's region, chosen from the schedule in advance.
\end{enumerate}
\textbf{Part IV.}
\begin{enumerate}[resume]
\item \emph{Named result}: in the model, a transaction from Tokyo reaches half the network in \qty{85.0}{\milli\second} and nine-tenths in
  \qty{99.5}{\milli\second} by flooding (113.9 and 149.0 with square-root fan-out); sending to the \omterm{def:nw:on-chain-connectivity:engine}{block engine} nearest the next leader rather than to a
  public node in Tokyo raises the share of leaders reached in time from 48\% to 61\% over a \qty{150}{\milli\second} window, and reaches a Frankfurt
  leader in 69.5 instead of \qty{83.5}{\milli\second}.
\item Documentation: the engines' cities and endpoints, the slot and leader rules; assumptions: the peer graph, the \omterm{def:nw:the-north-american-map:factor}{route factor}, the processing times.
\item Peers chosen among near nodes would shorten each hop but add hops across oceans; gossip's times would move, not its extra hops' cost.
\item Near the engines of the regions where most leaders are, with a fast route to the others; for this chain, Europe and one site in Asia.
\item A second path, with less variance, for what the engine does not carry, and delivery to builders and validators outside the engine's reach.
\item Send marked transactions through each path and measure the slot in which each lands, per leader region; compare with the model's arrivals.
\item In the on-chain rows: the engines per region, the chooser, the private relays and propagation networks, and the nodes the firm runs.
\item Because the next leader's location, known in advance, decides the fastest endpoint, and it changes every few slots.
\end{enumerate}
\end{solution}

\begin{solution}{iq:nw:on-chain-connectivity:1}
By gossip: each node forwards it, or announces it, to its peers on first receipt, validating it first. It is slow because the peers are random, so the
path zigzags, and every hop adds validation and queueing; the time grows with the number of hops and their worst links.
\iqlookfor{Peer-to-peer hops; processing per hop; announce and pull; measured tails.}
\end{solution}

\begin{solution}{iq:nw:on-chain-connectivity:2}
A service near the validators that takes transactions and bundles and forwards them directly to the block producer, often with an auction for inclusion.
There are several because producers are spread around the world and the last hop should be short wherever the next one is.
\iqlookfor{Direct forwarding; tips or auction; regional endpoints.}
\end{solution}

\begin{solution}{iq:nw:on-chain-connectivity:3}
Read the schedule for the coming slots, map each leader to a region (published or measured), and send each transaction to the endpoint nearest the leader
that will still be producing when it arrives; keep connections warm in every region, and measure landing rates per leader.
\iqlookfor{Advance schedule; mapping validators to places; handovers; measurement.}
\end{solution}

\begin{solution}{iq:nw:on-chain-connectivity:4}
When the node is on the trading path: state that must be read with the lowest latency, transactions the firm does not want to share with a provider,
or request rates beyond a provider's plans; a hosted endpoint is fine for reading history and for everything that is not a race.
\iqlookfor{Latency; limits; privacy; cost of operation.}
\end{solution}

\begin{solution}{iq:nw:on-chain-connectivity:5}
To keep it out of view of bots that would trade ahead of it or around it (frontrunning, sandwiching), and sometimes for a refund of the value it creates.
\iqlookfor{Visibility; MEV; private order flow.}
\end{solution}

\begin{solution}{iq:nw:on-chain-connectivity:6}
Connections to \omterm{def:nw:on-chain-connectivity:engine}{block engines} or builders in each leader region, a chooser driven by the \omterm{def:nw:on-chain-connectivity:leader}{leader schedule} with a margin at the handovers, private relays
where visibility matters, the firm's own nodes in the main regions for state, hosted nodes as back-up readers, per-region rate governors, and landing
rates measured per leader region and per path.
\iqlookfor{Schedule-driven routing; regional presence; private paths; measurement and limits.}
\end{solution}
